% GENERATED by tools/build_m05_code_sheet.py -- do not edit.
\begin{codebox}{Box 2}{Find communities with Leiden}
result = karate.community_leiden("modularity", n_iterations=-1)
leiden = result.membership
print(leiden)
\end{codebox}
\begin{linenotes}
\item[1] \texttt{community\_leiden} runs the Leiden algorithm. It looks for \emph{communities}: groups with many friendships inside and few between groups. \texttt{"modularity"} is the score it tries to raise: how much denser the groups are than in a random network. \texttt{n\_iterations=-1} means keep improving until nothing changes. What it returns is stored as \texttt{result}.
\item[2] \texttt{result.membership} is the answer as a plain list: entry $i$ is the group number of node $i$. We store it as \texttt{leiden}.
\item[3] Print the list. You should see 34 numbers. Two nodes with the same number are in the same group. The numbers are only labels, and Leiden uses chance, so your list may differ from your neighbour's.
\end{linenotes}
